
\section{Detectabilidad, resolución y probabilidad}
La resolución espacial distingue detalles; la detectabilidad depende además de contraste, ruido, orientación, procedimiento y decisión del inspector o sistema. Una indicación menor que una dimensión nominal del equipo puede detectarse en ciertas condiciones, y un defecto mayor puede pasar inadvertido en otras. No se transforma la resolución del equipo en un tamaño garantizado de grieta detectable.

La probabilidad de detección, POD, se estima con una población y un procedimiento de ensayo. Debe acompañarse de incertidumbre y del criterio estadístico usado. También existen falsas llamadas y errores de dimensionado. Un resultado negativo no prueba tamaño cero; se incorpora a un argumento de cobertura y capacidad compatible con el programa de inspección. Un valor $a_{90/95}$, cuando se usa, expresa una convención de probabilidad y confianza, no una certeza individual universal.
\section{Comparación por mecanismo físico}
Penetrantes requiere abertura a una superficie accesible y compatible; una capa que selle la abertura puede ocultar el defecto. Partículas magnéticas necesita material ferromagnético y orientación de campo pertinente. Corrientes inducidas explora perturbaciones electromagnéticas en conductores y está condicionada por frecuencia, geometría, propiedades y lift-off. La profundidad pelicular de P28 describe atenuación ideal, no una garantía de detección.

Ultrasonidos necesita acoplamiento y un recorrido interpretable; interfaces, grano y orientación afectan señal. Radiografía representa atenuación integrada y responde de modo distinto a discontinuidades volumétricas y planos mal orientados; una proyección puede superponer geometrías. Termografía observa respuesta térmica y depende de excitación, profundidad y propiedades. Emisión acústica escucha eventos durante una solicitación, sin proporcionar automáticamente una imagen completa de defectos inactivos.
\ilustra{ultrasonidos}{Modelo pulso--eco: el tiempo de vuelo incluye ida y vuelta. La posición de los ecos es ilustrativa y no constituye una calibración de equipo.}
\Needspace{.48\textheight}
\section{Ensayar una propiedad y decidir aceptación}
Tracción proporciona una curva dependiente de probeta, velocidad, temperatura y medida; dureza es una respuesta local bajo un método concreto. Las correlaciones entre dureza y resistencia son empíricas y no universales. Impacto caracteriza una respuesta dinámica de una geometría y condición de ensayo; su energía no es $K_{IC}$. Fatiga y fluencia añaden historia temporal y dispersión.

Un ensayo o END produce información; la aceptación compara esa información con criterios aplicables y considera cobertura, incertidumbre y configuración. Conservar método, equipo, revisión, zona, resultados y personal permite interpretar el informe años después. Una casilla «OK» sin contexto puede ser insuficiente para reconstruir qué se examinó.
\repaso{¿Por qué dos END que exploran la misma pieza pueden dar resultados diferentes sin que uno sea necesariamente incorrecto?}{Pueden ser sensibles a mecanismos, orientaciones y profundidades distintas, con diferentes coberturas y umbrales. Debe compararse la capacidad física de cada técnica y el procedimiento aplicado antes de atribuir una contradicción.}
